% =============================================================================
% 1. INTRODUCTION
% Target: 1.5 pages (~750 words)
% =============================================================================

\section{Introduction}
\label{sec:introduction}

% ---------- ¶1: Motivation ----------
Software development lifecycle (SDLC) data is now routinely
machine-readable across version control, issue trackers, code-review
systems, and continuous-integration platforms. Process mining (PM)
infers process structure from event logs; stochastic modelling
forecasts completion times and absorption probabilities; and a growing
body of work in agile practice has popularised Monte Carlo forecasting
at the project
level~\cite{vacanti2015, savage2009flaw, magennis2015forecasting}.
Yet the joint deployment of these techniques for \emph{predictive} use
in the SDLC remains rare: works tend to discover process models
without forecasting, forecast lead times without process awareness, or
apply machine learning to features that ignore the underlying process
structure.

% ---------- ¶2: Gap (from SLR) ----------
A recent systematic literature review across 5{,}783 deduplicated
records and 404 confirmed primary studies~\cite{almeida2026slr}
quantifies this fragmentation: only a \emph{single} primary study in
the corpus jointly satisfies the three dimensions PM, stochastic
modelling, and forecasting, and that study targets generic business
processes rather than the SDLC. The literature is therefore
descriptively rich but predictively fragmented; no end-to-end pipeline
specialised to the SDLC and formally specified --- with input/output
contracts, correctness guarantees, and consistency properties ---
appears in the surveyed corpus. The absence of such a pipeline is the
problem this paper addresses.

% ---------- ¶3: Contributions ----------
We introduce \textsc{PATHCAST}, a formally specified method that
integrates event log construction, process discovery, absorbing Markov
chain estimation, process-aware Monte Carlo simulation, and a
machine-learning enhancement component into a single forecasting
pipeline for the SDLC. The contributions of this paper are:
\begin{enumerate}
    \item \textbf{C1.} A formal pipeline specification with
    contract-based stage composition
    (Section~\ref{sec:architecture}, Table~\ref{tab:contracts}) and a
    correctness theorem under explicit preconditions
    (Theorem~\ref{thm:pipeline-correctness}).
    \item \textbf{C2.} An empirically grounded stochastic model that
    couples a Markov transition kernel with empirical sojourn-time
    distributions (Section~\ref{sec:stages}), preserving the
    multi-modal, heavy-tailed shapes observed in software process
    sojourns without parametric assumptions.
    \item \textbf{C3.} A novel family of features for delay prediction
    derived from PM and Markov analysis
    (Section~\ref{sec:ml-component}, Table~\ref{tab:features}),
    instantiating the trajectory-conditional expectations
    $\mathbb{E}[T \mid \psi]$ and $P(\textup{Done} \mid \psi)$ as
    predictor inputs that are by construction inaccessible to
    repository-mining baselines.
    \item \textbf{C4.} Formal guarantees: asymptotic statistical
    consistency (Theorem~\ref{thm:consistency}), end-to-end
    distributional convergence with vanishing calibration error
    (Theorem~\ref{thm:pipeline-convergence}), and an
    information-theoretic advantage over process-blind throughput-based
    Monte Carlo (Theorem~\ref{thm:information}).
    \item \textbf{C5.} A closed-form illustrative example
    (Section~\ref{sec:numerical-example}) that instantiates the
    pipeline on a six-state SDLC chain, derives the fundamental
    matrix and absorption probabilities analytically, and verifies
    agreement with Monte Carlo simulation within $1\%$ on every
    quantity for which a closed form exists.
\end{enumerate}

% ---------- ¶4: Scope ----------
The present paper specifies the method and its theoretical properties.
The empirical validation across 48 open-source repositories ---
two-stage experimental design (pipeline validation and forecasting
accuracy), comparison against four baselines under the continuous
ranked probability score, and statistical analysis across repositories
with Bonferroni correction --- is the subject of a companion paper. We
summarise the evaluation protocol in
Section~\ref{sec:evaluation-plan}; no empirical numbers from real
repositories appear here.

% ---------- ¶5: Roadmap ----------
Section~\ref{sec:background} reviews process mining, absorbing Markov
chains, stochastic forecasting in software engineering, and the
pipeline-fragmentation gap.
Section~\ref{sec:architecture} introduces the pipeline architecture
and the five inter-stage contracts.
Section~\ref{sec:stages} specifies the four sequential stages.
Section~\ref{sec:ml-component} specifies the parallel ML enhancement
component. Section~\ref{sec:properties} establishes the formal
properties, with proofs.
Section~\ref{sec:numerical-example} presents the closed-form
numerical example with Monte Carlo cross-validation.
Section~\ref{sec:comparison} compares \textsc{PATHCAST} with related
approaches along five theoretical dimensions.
Section~\ref{sec:evaluation-plan} outlines the planned empirical
validation. Section~\ref{sec:conclusion} concludes.
